% Standalone common lemma; copy/input once per rank manuscript.
% Requires EGODimPQ bibliography entry and standard amsthm/math notation.
\begin{lemma}[Cyclic-module divisibility]\label{lem:cyclic-module-divisibility}
Let $H\cong C_p$, with $p$ prime, be a group of invertibles in a
fusion category over $\mathbb C$. Suppose that simple objects $X,Y,Z$
satisfy $hX\cong X$, $Yh\cong Y$, and $hZ\cong Z\cong Zh$ for every
$h\in H$. If $\operatorname{FPdim}(Z)/p$ is not an algebraic integer,
then $p\mid N_{XY}^{Z}$.
\end{lemma}
\begin{proof}
The rank-one module category generated by $X$ trivializes the pointed
associator on $H$. Thus $A=\bigoplus_{h\in H}h$ admits its regular
connected separable algebra structure. Since
$H^2(C_p,\mathbb C^\times)=0$, the stabilizer cocycles can be
trivialized to give $X$ a left $A$-module structure and $Y$ a right
$A$-module structure. Hence $X\otimes Y$ is an $A$-bimodule.

The $A$-bimodules supported on the fixed simple $Z$ correspond to
projective representations of $C_p^2$: the projective-representation
dimension equals the underlying multiplicity of $Z$. A scalar
two-cocycle on $C_p^2$ has either zero or nondegenerate commutator
pairing. Its twisted group algebra is accordingly
$\mathbb C^{p^2}$ or $M_p(\mathbb C)$. The first possibility would
give a simple bimodule of FP dimension
$\operatorname{FPdim}(Z)/\operatorname{FPdim}(A)
=\operatorname{FPdim}(Z)/p$, contrary to algebraic integrality.
In the second possibility every simple bimodule supported on $Z$
has underlying object $pZ$. Decomposing the $Z$-isotypic bimodule
summand of $X\otimes Y$ proves the assertion. This is the regular
cyclic-algebra method of \cite[Lemmas~6.1--6.2]{EGODimPQ}, with
the stabilizer cocycle retained rather than assumed trivial.
\end{proof}
